\subsection*{Questão 12}

\textit{Se a largura de banda de um sinal composto vale 5 \khz e a menor frequência vale 52 \khz, qual o valor da maior frequência?}

\textbf{Resposta:}
A largura de banda ($B$) de um sinal composto é a diferença entre a sua maior e menor frequência:

\begin{equation}
    B = f_{max} - f_{min}
\end{equation}

Substituindo os valores fornecidos:
\begin{itemize}
    \item $B = 5 \khz$
    \item $f_{min} = 52 \khz$
\end{itemize}

\begin{equation}
    5 \khz = f_{max} - 52 \khz
    \implies f_{max} = 5 \khz + 52 \khz = 57 \khz
\end{equation}

\textbf{Resultado:}
A maior frequência do sinal composto é $57 \khz$.